\section{Clean Baseline}
\label{sec:clean-baseline}

\subsection{Purpose and Evidence Scope}

Before any Copy Fail-specific testing began, VM~270
\code{copyfail01} was recorded as the reference system. This was the state
against which later kernel, package, configuration, target-file, and privilege
changes could be compared. Preserving it separately also meant that the
vulnerable reconstruction on VM~271 did not overwrite the known starting
point.

Here, \emph{clean} describes a position in the experimental sequence. It means
that no vulnerability-specific configuration, proof-of-concept execution, or
exploitation had yet been introduced within the documented workflow. It does
not mean that every security property of the guest had already been tested, or
that the installed kernel could be classified as vulnerable or protected from
the baseline alone. Chapter~4 performs that narrower validation.

\begin{evidencebox}
The public guest-state evidence is stored under
\path{evidence/00-baseline/}. It contains 20 command-output artifacts, a
private-origin manifest, and a separate manifest for the sanitized public
copy. The reviewed Phase~00 screenshot directory contains no retained image,
so this chapter uses the machine-readable captures directly instead of
presenting a later-phase screenshot as baseline evidence.
\end{evidencebox}

The baseline covers the operating-system identity, active kernel, installed
packages, available upgrades, account state, network configuration, listening
sockets, and storage layout. A separate Proxmox transcript records the virtual
hardware and snapshot state. Together, the guest and hypervisor views establish
which machine was recorded and which checkpoint later experiments used as
their reference.

\subsection{Operating-System and Account State}

The guest identified itself as \code{copyfail01}, running Ubuntu 24.04.2 LTS
on x86-64 KVM/QEMU hardware. The relevant sanitized excerpt from
\path{command-output/hostnamectl.txt} is reproduced below.

\begin{lstlisting}
Static hostname: copyfail01
Chassis: vm
Virtualization: kvm
Operating System: Ubuntu 24.04.2 LTS
Kernel: Linux 6.8.0-137-generic
Architecture: x86-64
Hardware Vendor: QEMU
Hardware Model: Standard PC _Q35 + ICH9, 2009_
\end{lstlisting}

\Cref{tab:baseline-system-identity} combines that guest view with the
non-sensitive fields retained from the later hypervisor verification.

\begin{table}[H]
    \centering
    \caption{Recorded identity of the reference guest.}
    \label{tab:baseline-system-identity}
    \small
    \begin{tabularx}{0.92\textwidth}{@{}>{\bfseries}p{0.31\textwidth}Y@{}}
        \toprule
        Property & Recorded state \\
        \midrule
        Hypervisor object &
        Proxmox VM~270, configured as \code{copyfail01} \\

        Guest hostname &
        \code{copyfail01} \\

        Operating system &
        Ubuntu 24.04.2 LTS (Noble Numbat) \\

        Architecture &
        \code{x86-64} \\

        Virtual platform &
        KVM guest on QEMU Q35 hardware with UEFI firmware \\

        Administrative account &
        \code{researcher}, UID~1000 and GID~1000 \\
        \bottomrule
    \end{tabularx}
\end{table}

The matching hypervisor name and guest hostname link the command-output set to
the intended VM without depending on only one naming source. The
\code{id.txt} capture also records the baseline administrative account as
\code{researcher}, including membership in \code{adm}, \code{sudo}, and
\code{lxd}. This account was used for preparation and evidence collection; it
was not the unprivileged attacker identity used for the vulnerable
reproduction.

Machine and boot identifiers were redacted from the publication copy. Those
values are unnecessary for interpreting the experiment, while the hostname,
release, architecture, and virtualization fields retain the system identity
needed for comparison.

\subsection{Kernel and Package State}

Three runtime captures independently report the same active kernel release:
\path{kernel.txt}, \path{uname.txt}, and \path{proc-version.txt}. The package
database then links that active release to the installed Ubuntu image package.
This matters because an installed image does not prove which kernel is running,
and a runtime release string does not preserve the corresponding package
version.

\begin{lstlisting}
$ cat command-output/kernel.txt
6.8.0-137-generic

$ cat command-output/kernel-packages.txt
linux-image-6.8.0-137-generic    6.8.0-137.137
linux-image-generic              6.8.0-137.137

$ cat command-output/kernel-sha256.txt
76a7f2ef15fcbd2f5c25cd7e7b413f903b2078396063557f1dffb4a0b089a964  /boot/vmlinuz-6.8.0-137-generic
\end{lstlisting}

\begin{table}[H]
    \centering
    \caption{Recorded kernel state of VM~270.}
    \label{tab:baseline-kernel-state}
    \small
    \begin{tabularx}{\textwidth}{@{}>{\bfseries}p{0.29\textwidth}Y@{}}
        \toprule
        Property & Recorded state \\
        \midrule
        Running release &
        \code{6.8.0-137-generic} \\

        Kernel build &
        \code{\#137-Ubuntu SMP PREEMPT\_DYNAMIC}, built 17 July 2026 at
        20:28:23 UTC \\

        Installed image &
        \code{linux-image-6.8.0-137-generic}, version
        \code{6.8.0-137.137} \\

        Tracking package &
        \code{linux-image-generic}, version \code{6.8.0-137.137} \\

        Image architecture &
        \code{amd64} \\

        Recorded image digest &
        Captured in \path{command-output/kernel-sha256.txt} and reproduced
        in the evidence excerpt above \\
        \bottomrule
    \end{tabularx}
\end{table}

The digest identifies the kernel image that was present when the capture was
made. Because the image itself is not part of the public evidence set, the
digest cannot be recalculated from the repository alone. Manifest verification
confirms the integrity of the retained digest record, not the bytes of the
unpublished \code{vmlinuz} file.

The complete package inventory contains 680 package rows. The separate
\code{upgradable.txt} capture contains 77 package entries after excluding the
\code{apt} warning and heading lines. Those updates were recorded but not
applied before the baseline checkpoint.

\begin{table}[H]
    \centering
    \caption{Selected package facts relevant to later experimental work.}
    \label{tab:baseline-package-state}
    \small
    \begin{tabularx}{\textwidth}{@{}>{\bfseries\raggedright\arraybackslash}p{0.23\textwidth}>{\raggedright\arraybackslash}p{0.29\textwidth}Y@{}}
        \toprule
        Category & Recorded item & Baseline state \\
        \midrule
        Inventory scope &
        Installed packages &
        680 package rows with retained versions \\

        Pending changes &
        Upgradable packages &
        77 package rows; updates not applied before the checkpoint \\

        Kernel support &
        \code{kmod}; \code{initramfs-tools} &
        Versions \code{31+20240202-2ubuntu7.2} and
        \code{0.142ubuntu25.4} \\

        Existing observability &
        \code{bpfcc-tools}; \code{bpftrace}; \code{linux-tools-common} &
        Versions \code{0.29.1+ds-1ubuntu7},
        \code{0.20.2-1ubuntu4.3}, and \code{6.8.0-137.137} \\

        General tooling &
        \code{git}; \code{python3} &
        Versions \code{1:2.43.0-1ubuntu7.3} and
        \code{3.12.3-0ubuntu2.1} \\

        Exact build-tool entries &
        \code{build-essential}; \code{gcc}; \code{g++}; \code{clang};
        \code{make}; \code{bpftool}; \code{libbpf-dev} &
        No exact package row for these names in the captured inventory \\
        \bottomrule
    \end{tabularx}
\end{table}

The final row is intentionally narrow. Supporting components including
\code{gcc-14-base:amd64}, \code{libclang1-18}, and
\code{libclang-cpp18} were installed, so the evidence does not support the
broader claim that the system contained no GCC- or Clang-related packages. It
supports only the absence of the exact build-tool package names listed in the
table. Package presence also does not prove that a detector was executed, just
as package absence does not prove that no binary could have been transferred
to the guest.

\subsection{Network State}

The network baseline was recorded from both sides of the VM boundary. Proxmox
assigned VM~270 one VirtIO interface on \code{vmbr1} with firewall processing
enabled. Inside the guest, the corresponding interface appeared as
\code{enp6s18}; no second non-loopback interface appears in the captured address
state.

Netplan configured the interface for DHCP. The runtime route capture shows a
connected laboratory-subnet route and a default route through the laboratory
gateway. Concrete addresses and the MAC value were sanitized, while the
interface, addressing method, and routing role were retained.

\begin{lstlisting}
network:
  version: 2
  ethernets:
    enp6s18:
      dhcp4: true

default via <LAB_GATEWAY> dev enp6s18 proto dhcp src <LAB_VM_IP>
<LAB_SUBNET> dev enp6s18 proto kernel scope link src <LAB_VM_IP>
\end{lstlisting}

The listening-socket capture records SSH on TCP port~22 for IPv4 and IPv6,
plus the local \code{systemd-resolved} stub listeners on loopback port~53. No
other listener appears in that capture. This is a point-in-time result, not a
claim that the same listener set persisted through every later phase.

Together, the hypervisor and guest records match the topology from Chapter~2:
the Copy Fail guest used the isolated \code{vmbr1} segment and reached routed
destinations through the laboratory gateway rather than through a direct
management-network attachment.

\subsection{Storage State}

Proxmox presented a single 32~GiB virtual disk through VirtIO SCSI, with
discard, an I/O thread, and the virtual SSD flag enabled. The guest reported
the same capacity as \code{/dev/sda}. Its recorded layout is summarized in
\cref{tab:baseline-storage-state}.

\begin{table}[H]
    \centering
    \caption{Recorded storage layout and utilisation.}
    \label{tab:baseline-storage-state}
    \small
    \begin{tabularx}{\textwidth}{@{}>{\bfseries}p{0.27\textwidth}Y@{}}
        \toprule
        Component & Recorded state \\
        \midrule
        Virtual disk &
        One 32~GiB disk, visible inside the guest as \code{/dev/sda} \\

        EFI partition &
        \code{/dev/sda1}, 1~GiB VFAT, mounted at \code{/boot/efi} \\

        Boot partition &
        \code{/dev/sda2}, 2~GiB ext4, mounted at \code{/boot} \\

        LVM partition &
        \code{/dev/sda3}, 28.9~GiB, assigned to \code{ubuntu-vg} \\

        Root logical volume &
        \code{ubuntu-lv}, 14.47~GiB ext4, mounted at \code{/} \\

        Free LVM space &
        Approximately 14.47~GiB \\

        Root utilisation &
        15~GiB reported size, 5.6~GiB used, 7.9~GiB available,
        42\% utilisation \\
        \bottomrule
    \end{tabularx}
\end{table}

The storage capture establishes capacity and allocation before later
experiment-specific files were introduced. It does not imply that this layout
caused or prevented Copy Fail; it provides a fixed comparison point for later
file and filesystem observations.

\subsection{Integrity Verification}

The baseline uses separate SHA-256 manifests for two different byte sets. The
private manifest describes the original command output. The public manifest
describes the sanitized publication copy after machine and network identifiers
were removed. Reusing the private hashes for the modified public files would
have made legitimate redaction look like corruption, so the two states remain
explicitly separate.

An independent check of
\path{evidence/00-baseline/hashes/PUBLIC-SHA256SUMS} against the reviewed public
repository returned \code{OK} for all 20 enumerated files. The public command
output was also compared with the private-origin manifest to locate the effect
of sanitization. Seventeen files still matched; the expected mismatches were
\path{hostnamectl.txt}, \path{ip-addresses.txt}, and \path{routes.txt}.

\begin{table}[H]
    \centering
    \caption{Independent Phase~00 integrity results.}
    \label{tab:baseline-integrity-results}
    \small
    \begin{tabularx}{\textwidth}{@{}>{\bfseries\raggedright\arraybackslash}p{0.28\textwidth}>{\raggedright\arraybackslash}p{0.18\textwidth}Y@{}}
        \toprule
        Check & Result & Interpretation \\
        \midrule
        Public files against public manifest &
        20 of 20 matched &
        Every command-output object listed by the public manifest retained its
        recorded public bytes. \\

        Sanitized files against private-origin manifest &
        17 of 20 matched &
        The three mismatches correspond to the files changed by publication
        redaction. \\
        \bottomrule
    \end{tabularx}
\end{table}

The manifest scope matters. The 20 entries cover the command-output artifacts;
they do not cover the README, the manifests themselves, or every file beside
them. A successful check therefore detects changes to the listed bytes, but it
does not authenticate the origin of a command, prove that a capture was
complete, or settle a disagreement in documentation that lies outside the
manifest.

\subsection{Snapshot Verification and Baseline Result}

After the guest-state collection, Proxmox preserved VM~270 at the named
baseline checkpoint. Its exact snapshot name appears in the transcript below.
A later read-only capture, taken on 13 August 2026 at 21:53:07+02:00, found
the VM stopped and recorded both its parent field and its snapshot tree.

\begin{lstlisting}
status: stopped
cores: 2
memory: 4096
machine: q35
name: copyfail01
net0: virtio=<REDACTED_MAC>,bridge=vmbr1,firewall=1
parent: clean-baseline-00
scsi0: <REDACTED_STORAGE_ID>,discard=on,iothread=1,size=32G,ssd=1

`-> clean-baseline-00  2026-08-08 00:46:54
    Fresh Ubuntu 24.04.2 LTS Copy Fail research target.
 `-> current
\end{lstlisting}

The Phase~00 README calls the snapshot \code{00-clean-baseline}. That name does
not match the Proxmox evidence. Both the hypervisor parent field and the
snapshot listing use \code{clean-baseline-00}, so this chapter treats that as
the authoritative snapshot name. The README is not one of the 20 files covered
by the public command-output manifest, which is why the successful manifest
check does not resolve or conceal this documentation mismatch.

The snapshot listing records its creation time as 8 August 2026 at 00:46:54
and does not include a timezone. None is assigned here. The later transcript
does prove that the snapshot existed, that VM~270's then-current state descended
from it, and that the VM was stopped when the verification was collected. It
does not independently prove the content of an unexported virtual disk or every
historical action performed inside the guest.

At the end of this chapter, the experiment therefore has a concrete reference
state: Ubuntu 24.04.2 LTS on kernel \code{6.8.0-137-generic}, a documented
package and account inventory, one isolated laboratory interface, a recorded
storage layout, 20 publicly verifiable command-output artifacts, and a confirmed
Proxmox rollback point. That is enough to compare later state changes without
turning the word \emph{clean} into an unsupported security verdict.

Chapter~4 next checks \code{algif\_aead} availability, the effective loading
policy, and the kernel's patch state before VM~271 is modified.